\documentclass[12pt,a4paper]{article}
\usepackage[utf8]{inputenc}
\usepackage[T5]{fontenc}
\usepackage[vietnamese]{babel}
\usepackage{lmodern}
\usepackage[a4paper,margin=2.3cm]{geometry}
\usepackage{amsmath}
\usepackage{booktabs}
\usepackage{longtable}
\usepackage{array}
\usepackage{xcolor}
\usepackage{listings}
\usepackage[hidelinks]{hyperref}
\usepackage{enumitem}

\lstset{
  basicstyle=\ttfamily\small,
  breaklines=true,
  frame=single,
  columns=fullflexible,
  showstringspaces=false,
  keywordstyle=\color{blue!70!black},
  commentstyle=\color{green!40!black}
}

\newcommand{\code}[1]{\texttt{\detokenize{#1}}}

\title{Báo cáo thiết kế RTL CPU 4-bit}
\author{}
\date{\today}

\begin{document}
\maketitle
\tableofcontents
\newpage

\section{Mục tiêu thiết kế}

Thiết kế xây dựng một CPU đồng bộ có datapath 4-bit để minh họa đầy đủ
chu trình thiết kế số từ RTL đến physical design. CPU thực thi instruction
8-bit, hỗ trợ các phép toán số học, logic, dịch bit, nhảy và dừng. RTL
được viết bằng SystemVerilog và có thể tổng hợp bằng Yosys sang standard
cell Sky130.

Các mục tiêu chính:

\begin{itemize}[noitemsep]
  \item datapath và địa chỉ rộng 4 bit;
  \item instruction rộng 8 bit;
  \item kiến trúc accumulator đơn giản;
  \item bộ điều khiển FSM ba trạng thái;
  \item reset bất đồng bộ active-low;
  \item instruction memory nằm ngoài CPU;
  \item RTL synthesizable, không suy diễn latch;
  \item có self-checking testbench và chương trình kiểm thử.
\end{itemize}

\section{Cấu trúc source code}

\begin{longtable}{>{\raggedright\arraybackslash}p{0.34\textwidth}
                  >{\raggedright\arraybackslash}p{0.56\textwidth}}
\toprule
\textbf{File} & \textbf{Chức năng} \\
\midrule
\endhead
\code{cpu_4bit.sv} & Top-level CPU, các thanh ghi trạng thái và kết nối datapath. \\
\code{control_unit.sv} & Giải mã instruction và điều khiển FSM. \\
\code{alu_4bit.sv} & Thực hiện các phép toán số học/logic và tạo flags. \\
\code{program_rom.sv} & Mô hình ROM 16 từ dùng cho mô phỏng độc lập. \\
\code{tb.sv} & Self-checking RTL testbench. \\
\code{memory/program.hex} & Chương trình kiểm thử 16 byte. \\
\code{filelist.f} & Danh sách source dùng cho Xcelium. \\
\bottomrule
\end{longtable}

\section{Giao diện top-level}

Module \code{cpu_4bit} có giao diện:

\begin{lstlisting}[language=Verilog]
module cpu_4bit (
    input  logic       clk,
    input  logic       rst_n,
    output logic [3:0] imem_addr,
    input  logic [7:0] imem_rdata,
    output logic       halted
);
\end{lstlisting}

\begin{longtable}{>{\raggedright\arraybackslash}p{0.20\textwidth}
                  >{\centering\arraybackslash}p{0.12\textwidth}
                  >{\raggedright\arraybackslash}p{0.54\textwidth}}
\toprule
\textbf{Signal} & \textbf{Width} & \textbf{Ý nghĩa} \\
\midrule
\endhead
\code{clk} & 1 & Clock cạnh lên. Physical design và testbench cuối dùng chu kỳ 20\,ns. \\
\code{rst_n} & 1 & Reset bất đồng bộ, tác động mức thấp. \\
\code{imem_addr} & 4 & Địa chỉ instruction, lấy trực tiếp từ program counter. \\
\code{imem_rdata} & 8 & Instruction đọc bất đồng bộ từ bộ nhớ ngoài. \\
\code{halted} & 1 & Bằng 1 khi CPU đã thực thi lệnh HLT. \\
\bottomrule
\end{longtable}

Instruction memory không nằm trong core. Cách tách này giúp CPU có giao
diện rõ ràng, tránh tổng hợp ROM test thành standard cells và cho phép
thay đổi chương trình mà không sửa datapath.

\section{Kiến trúc datapath}

\subsection{Các thanh ghi}

CPU chứa các thanh ghi sau:

\begin{itemize}[noitemsep]
  \item \code{pc[3:0]}: program counter, tạo địa chỉ 0--15;
  \item \code{ir[7:0]}: instruction register;
  \item \code{acc[3:0]}: accumulator và toán hạng chính của ALU;
  \item \code{zero_flag}, \code{carry_flag}, \code{borrow_flag},
        \code{negative_flag}, \code{overflow_flag}: các cờ trạng thái.
\end{itemize}

Instruction được chia thành:

\[
\text{instruction}[7:0]
= \underbrace{\text{opcode}[3:0]}_{\text{IR}[7:4]}
\parallel
\underbrace{\text{operand}[3:0]}_{\text{IR}[3:0]}.
\]

Opcode chọn phép toán, còn operand có thể là immediate data hoặc jump
target. Đây là encoding phù hợp với CPU nhỏ vì mỗi instruction chỉ cần
một byte.

\subsection{Cập nhật tuần tự}

Các thanh ghi dùng nonblocking assignment:

\begin{lstlisting}[language=Verilog]
always_ff @(posedge clk or negedge rst_n) begin
  if (!rst_n) begin
    pc  <= 4'b0000;
    ir  <= 8'b0000_0000;
    acc <= 4'b0000;
  end else begin
    if (ir_load)  ir  <= imem_rdata;
    if (pc_load)  pc  <= operand;
    else if (pc_inc) pc <= pc + 4'd1;
    if (acc_write) acc <= alu_result;
  end
end
\end{lstlisting}

Toán tử \code{<=} cập nhật thanh ghi ở NBA region ngay sau cạnh clock
hiện tại; nó không tự tạo thêm một chu kỳ. Một chu kỳ pipeline chỉ xuất
hiện khi dữ liệu phải đi qua hai thanh ghi liên tiếp hoặc khi FSM cố ý
chia thao tác thành nhiều trạng thái.

\section{Bộ điều khiển FSM}

Control unit có ba trạng thái:

\begin{enumerate}[noitemsep]
  \item \textbf{FETCH}: nạp \code{IR <- imem_rdata}, tăng
        \code{PC <- PC + 1}, sau đó chuyển sang EXECUTE;
  \item \textbf{EXECUTE}: giải mã opcode, phát control signals và thực thi
        instruction, thông thường quay lại FETCH;
  \item \textbf{HALT}: giữ \code{halted=1} và không thực thi thêm lệnh.
\end{enumerate}

Luồng trạng thái cơ bản:

\[
\text{FETCH} \rightarrow \text{EXECUTE}
\rightarrow \text{FETCH} \rightarrow \cdots
\rightarrow \text{HALT}.
\]

Vì mỗi instruction dùng một chu kỳ FETCH và một chu kỳ EXECUTE, throughput
cơ bản là một instruction trong hai chu kỳ. Đây là nguyên nhân
\code{imem_addr} giữ nguyên qua hai clock phases; không phải do
nonblocking assignment.

\section{Tập lệnh}

\begin{longtable}{>{\centering\arraybackslash}p{0.11\textwidth}
                  >{\raggedright\arraybackslash}p{0.14\textwidth}
                  >{\raggedright\arraybackslash}p{0.58\textwidth}}
\toprule
\textbf{Opcode} & \textbf{Mnemonic} & \textbf{Hoạt động} \\
\midrule
\endhead
\code{0x0} & NOP & Không thay đổi trạng thái datapath. \\
\code{0x1} & LDI & \code{ACC <- operand}. \\
\code{0x2} & ADD & \code{ACC <- ACC + operand}. \\
\code{0x3} & SUB & \code{ACC <- ACC - operand}. \\
\code{0x4} & AND & \code{ACC <- ACC & operand}. \\
\code{0x5} & OR  & \code{ACC <- ACC | operand}. \\
\code{0x6} & XOR & \code{ACC <- ACC ^ operand}. \\
\code{0x7} & NOT & \code{ACC <- ~ACC}. \\
\code{0x8} & SHL & Dịch trái ACC một bit. \\
\code{0x9} & SHR & Dịch phải logic ACC một bit. \\
\code{0xA} & JMP & \code{PC <- operand}. \\
\code{0xB} & JZ  & Nạp PC bằng operand nếu zero flag bằng 1. \\
\code{0xC} & INC & \code{ACC <- ACC + 1}. \\
\code{0xD} & DEC & \code{ACC <- ACC - 1}. \\
\code{0xF} & HLT & Chuyển FSM sang HALT. \\
\bottomrule
\end{longtable}

Opcode \code{0xE} chưa được định nghĩa; control unit xử lý opcode không
hợp lệ giống NOP để thiết kế luôn có hành vi xác định.

\section{ALU 4-bit và flags}

ALU là logic tổ hợp, nhận hai toán hạng \code{a}, \code{b} cùng
\code{alu_op}. Mọi output được gán giá trị mặc định trước câu lệnh
\code{case}, do đó không suy diễn latch.

\begin{longtable}{>{\raggedright\arraybackslash}p{0.20\textwidth}
                  >{\raggedright\arraybackslash}p{0.66\textwidth}}
\toprule
\textbf{Flag} & \textbf{Ý nghĩa} \\
\midrule
\endhead
\code{zero} & Result bằng 0. \\
\code{negative} & Bằng bit MSB của result, biểu diễn dấu trong two's complement. \\
\code{carry} & Carry-out của ADD/INC hoặc bit bị dịch ra trong SHL/SHR. \\
\code{borrow} & Bằng 1 khi phép trừ unsigned có \code{a < b}; DEC tại 0 cũng tạo borrow. \\
\code{overflow} & Tràn signed của ADD, SUB, INC hoặc DEC. \\
\bottomrule
\end{longtable}

Ví dụ carry và overflow của phép cộng được tính độc lập:

\begin{lstlisting}[language=Verilog]
temp     = {1'b0, a} + {1'b0, b};
result   = temp[3:0];
carry    = temp[4];
overflow = (~(a[3] ^ b[3])) & (result[3] ^ a[3]);
\end{lstlisting}

Carry mô tả tràn unsigned, trong khi overflow mô tả tràn signed.

\section{Reset và coding style}

Reset \code{rst_n} là asynchronous active-low:

\begin{lstlisting}[language=Verilog]
always_ff @(posedge clk or negedge rst_n)
\end{lstlisting}

Khi reset, PC, IR, ACC, flags và FSM state đều về giá trị xác định. Logic
tuần tự dùng \code{always_ff} và nonblocking assignment; logic tổ hợp dùng
\code{always_comb} và blocking assignment. Yosys xác nhận không suy diễn
latch cho ALU hoặc control unit.

\section{Môi trường kiểm thử RTL}

Testbench tạo clock 20\,ns để khớp constraint physical design, giữ reset
qua hai cạnh lên rồi nhả reset. Instruction memory được mô hình hóa như
ROM bất đồng bộ zero-wait-state:

\begin{lstlisting}[language=Verilog]
assign imem_rdata = program_memory[imem_addr];
\end{lstlisting}

Chương trình kiểm thử:

\begin{longtable}{>{\centering\arraybackslash}p{0.12\textwidth}
                  >{\centering\arraybackslash}p{0.12\textwidth}
                  >{\raggedright\arraybackslash}p{0.24\textwidth}
                  >{\raggedright\arraybackslash}p{0.30\textwidth}}
\toprule
\textbf{Addr} & \textbf{Byte} & \textbf{Instruction} & \textbf{ACC sau execute} \\
\midrule
\endhead
0 & \code{13} & LDI 3 & 3 \\
1 & \code{25} & ADD 5 & 8 \\
2 & \code{32} & SUB 2 & 6 \\
3 & \code{80} & SHL & 12 \\
4 & \code{90} & SHR & 6 \\
5 & \code{F0} & HLT & 6 \\
\bottomrule
\end{longtable}

Testbench chờ \code{halted}, đợi thêm một cạnh clock rồi kiểm tra
\code{dut.acc == 4'd6}. Nếu CPU không halt trước 500\,ns, watchdog báo
fatal error.

\section{Chạy mô phỏng và xem waveform}

Lệnh Xcelium:

\begin{lstlisting}
xrun -64bit -sv -timescale 1ns/1ps -access +rwc -xmlibdirname runs/manual_tcl_util65/xcelium_rtl_20ns.d -f filelist.f -top cpu_4bit_tb -input xcelium_rtl_shm.tcl
\end{lstlisting}

Mở waveform:

\begin{lstlisting}
simvision runs/manual_tcl_util65/rtl_wave.shm
\end{lstlisting}

Waveform đúng có chuỗi địa chỉ
\code{0,1,2,3,4,5,6}, dữ liệu
\code{13,25,32,80,90,F0,00} và \code{halted} lên 1 sau khi HLT được
execute. RTL và post-layout waveform có cùng trình tự chức năng. RTL đổi
tín hiệu qua delta-cycle gần cạnh clock; post-layout có propagation delay
do SDF.

\section{Kết quả synthesis tham chiếu}

Run Yosys tại \code{runs/manual_pd_01/06-yosys-synthesis} cho kết quả:

\begin{longtable}{>{\raggedright\arraybackslash}p{0.50\textwidth}
                  >{\raggedleft\arraybackslash}p{0.30\textwidth}}
\toprule
\textbf{Chỉ tiêu} & \textbf{Kết quả} \\
\midrule
\endhead
Top module & \code{cpu_4bit} \\
Số wire sau technology mapping & 140 \\
Số standard-cell instance & 140 \\
Số flip-flop \code{dfrtp_2} & 19 \\
Mapped cell area & 1591.5264\,\(\mu\mathrm{m}^2\) \\
\bottomrule
\end{longtable}

Mười chín flip-flop tương ứng với PC 4 bit, IR 8 bit, ACC 4 bit, năm
flags và FSM state sau tối ưu/mapping. Synthesis tạo netlist standard-cell
làm đầu vào cho floorplan và các bước physical design.

\section{Giới hạn và hướng phát triển}

\begin{itemize}
  \item CPU chưa có data memory, stack hoặc register file.
  \item Mỗi instruction cần hai chu kỳ; có thể cải thiện bằng pipeline
        hoặc kiến trúc single-cycle.
  \item Instruction memory giả định asynchronous zero-wait-state; giao
        diện thực tế nên có handshake hoặc wait-state.
  \item JZ chỉ dùng zero flag; có thể bổ sung branch theo carry, negative
        và overflow.
  \item Cần testbench bổ sung để kiểm tra từng opcode, boundary values,
        signed overflow, carry và borrow.
  \item Có thể bổ sung assertions và functional coverage cho FSM,
        instruction decode và flags.
\end{itemize}

\section{Kết luận}

RTL CPU 4-bit có kiến trúc accumulator nhỏ gọn nhưng chứa đủ các thành
phần cơ bản của một bộ xử lý: PC, IR, ALU, flags, control FSM và giao tiếp
instruction memory. Coding style phân tách rõ logic tuần tự và tổ hợp,
không suy diễn latch và tổng hợp thành công sang Sky130. Chương trình kiểm
thử thực hiện đúng chuỗi lệnh và kết thúc với ACC bằng 6. Đây là nền tảng
để thực hiện synthesis, physical design, signoff và post-layout
gate-level simulation.

\end{document}
